\newpage
\addcontentsline{toc}{section}{Tài liệu tham khảo}


\begin{thebibliography}{10}
	
	\bibitem{cormen2022}
	Cormen, T. H., Leiserson, C. E., Rivest, R. L., \& Stein, C. (2022). \textit{Introduction to algorithms} (4th ed.). MIT Press.
	
	\bibitem{dangthienbinh}
	Đặng Thiên Bình (n.d.). \textit{Bài giảng Cấu trúc dữ liệu}. Lưu hành nội bộ.
	
	\bibitem{chahelgupta}
	chahelgupta. (n.d.). \textit{DSA-Project-FamilyTree: C++ program for building and visualizing a hierarchical family tree}. GitHub. \url{https://github.com/chahelgupta/DSA-Project-FamilyTree}
	
	\bibitem{praveenmylavarapu}
	praveenmylavarapu. (n.d.). \textit{family-tree: A family tree program in C++ using n-ary tree}. GitHub. \url{https://github.com/praveenmylavarapu/family-tree}
	
	\bibitem{tenouk}
	tenouk. (n.d.). \textit{Constructing a very simple family structure tree using the C++ inheritance principles}. tenouk.com. \url{https://www.tenouk.com/cpluscodesnippet/cplusclassinheritancefamilytree.html}
	
	\bibitem{qnu2025}
	Trường Đại học Quy Nhơn. (2025). \textit{Mẫu code C++ cho cây nhị phân và cây gia phả}. Studocu. \url{https://www.studocu.vn/vn/document/dai-hoc-quy-nhon/li-luan-giang-day/mau-code-c-cho-cay-nhi-phan-va-cay-gia-pha-kcj/148579038}
	
	\bibitem{minhtc}
	homielab. (n.d.). \textit{Gia Phả OS (Gia Phả Open Source)}. Github. \url{https://github.com/homielab/giapha-os}
	
	\bibitem{wikipedia2025}
	Wikipedia contributors. (2025). \textit{Kinship}. Wikipedia, The Free Encyclopedia. \url{https://en.wikipedia.org/wiki/Kinship}
	
\end{thebibliography}